% Stefan Zweig: The World of Yesterday: Reading Summary
% Introduction to the European Union, Week 1
% Compile with: pdflatex zweig-world-of-yesterday-reading-summary.tex

\documentclass[11pt,a4paper]{article}
\usepackage[utf8]{inputenc}
\usepackage[T1]{fontenc}
\usepackage[margin=2.5cm]{geometry}
\usepackage{booktabs}
\usepackage{enumitem}
\setlist{nosep,leftmargin=*}
\usepackage{parskip}
\usepackage{microtype}
\usepackage{hyperref}
\usepackage{xcolor}
\definecolor{eublue}{RGB}{0,51,153}

\hyphenation{Spitzenkandidat Spitzenkandidaten}
\hypersetup{
  colorlinks=true,
  linkcolor=eublue,
  urlcolor=eublue,
  pdftitle={Zweig: The World of Yesterday},
  pdfauthor={Introduction to the EU, UCM}
}

\begin{document}

\title{Stefan Zweig: \textit{The World of Yesterday}\\Key Quotes \& Themes --- Rhetorical Analysis}
\author{Introduction to the European Union\\Universidad Complutense de Madrid\\Week 1}
\date{}
\maketitle
\vspace{-1em}
\hrule
\vspace{1em}

\section*{Rhetorical Framework (Ethos, Logos, Pathos)}

\begin{table}[h]
\centering
\begin{tabular}{@{}lll@{}}
\toprule
\textbf{Appeal} & \textbf{Meaning} & \textbf{Zweig's Use} \\
\midrule
\textbf{Ethos} & Credibility, authority, trust & Personal witness; lived experience as refugee and exile \\
\textbf{Logos} & Logic, structure, argument & Historical analysis; cause-and-effect; structural parallels \\
\textbf{Pathos} & Emotion, stakes, persuasion & Nostalgia, grief, loss; appeals to reader's sense of belonging \\
\bottomrule
\end{tabular}
\end{table}

\textbf{Key insight:} Zweig blends all three. His \textit{ethos} comes from having lived through the collapse; his \textit{logos} from tracing cause and effect; his \textit{pathos} from the emotional weight of displacement and loss.

\section*{Bibliographic Information}

\begin{itemize}
\item \textbf{Author:} Stefan Zweig
\item \textbf{Title:} \textit{The World of Yesterday}
\item \textbf{Born:} 1881 in Habsburg Monarchy (Vienna)
\item \textbf{Died:} 1942 (suicide in exile)
\item \textbf{Context:} Jewish-Austrian writer who witnessed both World Wars; memoirs written in exile, published posthumously
\end{itemize}

\textbf{Ethos:} Zweig's authority rests on his position as first-hand witness---someone who experienced the ``before'' and ``after'' of Europe's collapse.

\section*{Summary \& Key Points}

\subsection*{I. Displacement and Loss of Home}

\textbf{Quote:}
\begin{quote}
``For I have indeed been torn from all my roots, even from the earth that nourished them, more entirely than most in our times.''

``So I belong nowhere now, I am a stranger or at the most a guest everywhere.''
\end{quote}

\textbf{Rhetorical analysis:}
\begin{itemize}
\item \textbf{Pathos:} Strong emotional appeal---loss of roots, belonging, identity. Reader feels the weight of exile.
\item \textbf{Ethos:} Personal testimony grounds the abstract idea of displacement in lived experience.
\item \textbf{Logos:} Implicit argument---uprooting destroys identity; without roots, one becomes a perpetual guest.
\end{itemize}

\textbf{Analysis:} Zweig's experience of losing his homeland, language, and identity reflects broader European displacement after the wars.

\subsection*{II. The Golden Age of Security}

\textbf{Quote:}
\begin{quote}
``IF I TRY TO FIND some useful phrase to sum up the time of my childhood and youth before the First World War, I hope I can put it most succinctly by calling it the Golden Age of Security.''
\end{quote}

\textbf{Rhetorical analysis:}
\begin{itemize}
\item \textbf{Pathos:} Nostalgia---idealizes a past that readers may recognize or long for.
\item \textbf{Logos:} Defines a period by a single characteristic (security); sets up contrast with what follows.
\item \textbf{Ethos:} Personal memory as evidence; ``I was there'' gives the claim weight.
\end{itemize}

\textbf{Analysis:} Nostalgia for pre-war Europe characterized by stability, open borders, and sense of progress.

\subsection*{III. Rise of Nationalism and Xenophobia}

\textbf{Quote:}
\begin{quote}
``Nationalism emerged to agitate the world only after the war, and the first visible phenomenon which this intellectual epidemic of our century brought about was xenophobia; morbid dislike of the foreigner, or at least fear of the foreigner.''
\end{quote}

\textbf{Rhetorical analysis:}
\begin{itemize}
\item \textbf{Logos:} Causal argument---nationalism $\to$ xenophobia; ``epidemic'' metaphor suggests spread and contagion.
\item \textbf{Pathos:} ``Morbid,'' ``fear of the foreigner''---evokes unease and danger.
\item \textbf{Ethos:} Observer's voice; diagnosis of a collective pathology.
\end{itemize}

\textbf{Analysis:} Connection between nationalism and fear of foreigners; ``intellectual epidemic'' spreads like disease.

\subsection*{IV. Bureaucratic Restrictions}

\textbf{Quote:}
\begin{quote}
``All the humiliations previously devised solely for criminals were now inflicted on every traveller before and during every journey.''
\end{quote}

\textbf{Rhetorical analysis:}
\begin{itemize}
\item \textbf{Pathos:} ``Humiliations''---emotional charge; ordinary people treated like criminals.
\item \textbf{Logos:} Implicit comparison---before: citizens; after: suspects. Structural change in borders.
\item \textbf{Ethos:} Personal experience of crossing borders; witness to degradation.
\end{itemize}

\textbf{Analysis:} Transformation of borders from ``symbolic lines'' to barriers with extensive documentation requirements.

\subsection*{V. The Shattering Moment}

\textbf{Quote:}
\begin{quote}
``Then, on 28th June 1914, a shot was fired in Sarajevo, the shot that in a single second was to shatter the world of security and creative reason in which we had been reared\ldots''
\end{quote}

\textbf{Rhetorical analysis:}
\begin{itemize}
\item \textbf{Pathos:} ``Shatter,'' ``hollow clay pot,'' ``thousand pieces''---violence of sudden destruction.
\item \textbf{Logos:} Single event (assassination) as cause; before/after structure; ``one second'' = turning point.
\item \textbf{Ethos:} ``We had been reared''---collective ``we'' includes reader; shared experience of loss.
\end{itemize}

\textbf{Analysis:} Assassination of Archduke Franz Ferdinand as turning point; sudden destruction of old world.

\subsection*{VI. European Identity}

\textbf{Quote:}
\begin{quote}
``Nowhere was it easier to be a European, and I know that in part I have to thank Vienna\ldots for the fact that I learnt early to love the idea of community as the highest ideal of my heart.''
\end{quote}

\textbf{Rhetorical analysis:}
\begin{itemize}
\item \textbf{Ethos:} Vienna as cultural capital; lineage to Roman values; personal formation.
\item \textbf{Pathos:} ``Highest ideal of my heart''---emotional commitment to community.
\item \textbf{Logos:} Vienna = model of European cosmopolitanism; community as political ideal.
\end{itemize}

\section*{Main Arguments}

\begin{itemize}
\item Displacement and loss of home destroy identity; without roots one becomes a perpetual guest.
\item Pre-war Europe was a ``Golden Age of Security''---contrasted with postwar chaos.
\item Nationalism emerged as an ``intellectual epidemic'' leading to xenophobia.
\item Borders transformed from symbolic lines to barriers with humiliating documentation.
\item The Sarajevo assassination shattered the world of security in a single moment.
\item Vienna exemplified European cosmopolitanism; community as highest ideal.
\end{itemize}

\section*{Context \& Analysis}

\textbf{Debate Topics from Zweig (Rhetorical Framing):}

\begin{table}[h]
\centering
\footnotesize
\begin{tabular}{@{}llll@{}}
\toprule
\textbf{Topic} & \textbf{Ethos} & \textbf{Logos} & \textbf{Pathos} \\
\midrule
Personal vs.\ objective history & Whose testimony? & What evidence? & Whose story moves us? \\
Progress: declined or changed? & Who defines progress? & What metrics? & What do we long for? \\
Internationalism vs.\ nationalism & Who speaks for ``the people''? & Mutually exclusive? & Fear vs.\ belonging \\
Collective memory & Who shapes it? & How does it change? & What do we forget? \\
\bottomrule
\end{tabular}
\end{table}

\section*{Relevance to Course}

\begin{itemize}
\item \textbf{Ethos:} Understanding why European integration was necessary---who had lived through war and wanted peace.
\item \textbf{Logos:} Context for post-WWII desire for peace and cooperation; cause-and-effect of nationalism.
\item \textbf{Pathos:} Warning about dangers of nationalism and fear-based politics; emotional stakes of ``Europe.''
\item \textbf{All three:} Reflection on what ``Europe'' means as an identity---credibility, logic, and feeling together.
\end{itemize}

\vspace{1em}
\noindent\footnotesize\textit{Introduction to the European Union, Universidad Complutense de Madrid.}

\end{document}
